\section{Measure wrap work on a fixed operation}\label{ring-measurement}
The verified forward and reverse copy implementations satisfy the same contract.
They use different per-byte address calculations. Compare complete successful calls
with capacity eight, rather than timing an entire IDE launch or including assembly.
For push, the starting index is tail in an empty queue. For pop, it is head in a
queue containing one record. Each row below describes that initial state.

\Needspace{15\baselineskip}
\begin{longtable}{lrrrr}
  \toprule Call & Payload bytes & Start index & Forward instr. & Reverse instr. \\\midrule
  Push          & 0             & 7           & 32             & 32             \\
  Push          & 3             & 0           & 61             & 73             \\
  Push          & 3             & 7           & 62             & 77             \\
  Push          & 7             & 4           & 102            & 134            \\
  Pop           & 0             & 7           & 30             & 30             \\
  Pop           & 3             & 0           & 59             & 71             \\
  Pop           & 3             & 7           & 60             & 75             \\
  Pop           & 7             & 4           & 100            & 132            \\\bottomrule
\end{longtable}
The counter observes native completed forward instructions from the function's first
instruction through its return, stopping before the supplied continuation. Caller
argument setup, session creation, assembly, main, and exit are excluded. The inputs
are valid and the destination is large enough for each successful pop. These are
counts for the two measured implementations, not requirements for student source.

The forward loop retains a wrapped index between bytes. Moving the three-byte
record start from zero to seven adds one executed reset instruction. The reverse
loop computes each byte's circular index from its logical offset, then reduces it.
All three payload indices cross the boundary in the latter example, and its final
state calculation also needs a subtraction. Its count rises by four instructions.
The reverse implementation is useful evidence that grading accepts another legal
copy order; it is not an optimization over this forward loop.

For these successful forward paths, each payload byte adds ten instructions, and
the complete record crosses the physical end at most once. Counting the fixed work
and the optional reset explains \(31+10L+W\) for push and \(29+10L+W\) for pop,
where \(W\) is one if the record advance reaches or crosses the end and zero
otherwise. A zero-length record at the last slot still has \(W=1\), because its
prefix advances the cursor.

A two-span copy can move wrap decisions outside an inner byte loop. Before claiming
a saving, derive the span sizes, handle empty payloads, and test a prefix at the
last slot. The exercise checks permit that implementation. Instruction counts alone
do not establish host elapsed time or cycles on a physical MIPS processor. Measure
the relevant quantity separately and retain the same memory and failure contract.
